% 05-odd.tex -- the odd family.  Formulas F12-F13.
% Input by mostlyharmless.tex.  No preamble here.

\section{The odd family}
\label{sec:odd}

This section evaluates $X_{lm}(q)$ of \eqref{eq:radial-integral} for $l+m$ odd, and
does so completely: the answer is a single spherical Bessel function.  The even
family is section 6.

\subsection{The parity of the integrand}
\label{sec:odd-parity}

Write $\rho = \sin\theta$, so that $\sqrt{1-\rho^2} = \cos\theta$, and the radial
integral becomes
\begin{equation}
  X_{lm}(q) = \int_0^{\pi/2} P_{lm}(\cos\theta)\,
  J_m(2\pi q\sin\theta)\,\sin\theta\cos\theta\,d\theta .
  \label{eq:radial-in-theta}
\end{equation}
For $l+m$ odd the associated Legendre function is odd in its argument,
$P_{lm}(-z) = -P_{lm}(z)$, and so is the explicit factor $\cos\theta$; everything
else in the integrand is even.  The product is therefore even and the half-range
integral can be doubled,
\begin{equation}
  2X_{lm}(q) = \int_{-1}^{1} P_{lm}(z)\, J_m\!\left(2\pi q\sqrt{1-z^2}\right)
  z\,dz , \qquad z = \cos\theta ,
  \label{eq:doubled}
\end{equation}
which is the step the odd family turns on and which is available only here.

\subsection{Absorbing the factor of \texorpdfstring{$z$}{z}}
\label{sec:odd-recurrence}

The Legendre recurrence \citep[12.92]{arfken} removes that factor,
\begin{equation}
  z P_{lm}(z) = \frac{l-m+1}{2l+1} P_{l+1}^m(z)
  + \frac{l+m}{2l+1} P_{l-1}^m(z) ,
  \label{eq:legendre-recurrence}
\end{equation}
leaving two integrals of one form,
\begin{equation}
  K_{lm}(a) = \int_{-1}^{1} P_{lm}(z)\, J_m\!\left(a\sqrt{1-z^2}\right) dz ,
  \qquad a = 2\pi q ,
  \label{eq:K}
\end{equation}
so that
\begin{equation}
  2X_{lm}(q) = \frac{l-m+1}{2l+1} K_{l+1}^m
  + \frac{l+m}{2l+1} K_{l-1}^m .
  \label{eq:K-combination}
\end{equation}

\subsection{The \texorpdfstring{$K$}{K} identity}
\label{sec:k-identity}

For $l+m$ odd, \citet{neves2006} give
\begin{equation}
  K_{lm}(a) = (-1)^m \, 2 \, \frac{(l+m-1)!!}{(l-m)!!} \, j_l(a) .
  \label{eq:K-identity}
\end{equation}
This is the one step we cite rather than derive, and the overall sign is the one thing the
citation does not settle for us: it depends on the phase convention a source carries inside
$P_{lm}$, which is not always the one used here.  Certificate \vtag{18} therefore checks
\eqref{eq:K-identity} numerically in our own convention rather than taking the sign on
trust.  \emph{Numerical}, $l \le 6$.

\subsection{The collapse}
\label{sec:odd-collapse}

Substituting \eqref{eq:K-identity} into \eqref{eq:K-combination} gives two spherical
Bessel functions of different order, with coefficients that look different.  They
are not:
\begin{equation}
  \frac{l-m+1}{2l+1}\frac{(l+m)!!}{(l-m+1)!!}
  = \frac{l+m}{2l+1}\frac{(l+m-2)!!}{(l-m-1)!!}
  = \frac{(l+m)!!}{(2l+1)(l-m-1)!!} ,
  \label{eq:coefficients-coincide}
\end{equation}
using $(l-m+1)!! = (l-m+1)(l-m-1)!!$ on the left and
$(l+m)(l+m-2)!! = (l+m)!!$ on the right.  The middle expression is the one that
takes care: the second $K$ is $K_{l-1}^m$, so its double factorial starts at
$l+m-2$, not $l+m-1$.  Both terms
carry the same coefficient, so two adjacent Bessel orders stand side by side and the
recurrence \citep[10.51.1]{dlmf}
\begin{equation}
  j_{l-1}(a) + j_{l+1}(a) = \frac{2l+1}{a} j_l(a)
  \label{eq:bessel-recurrence}
\end{equation}
fuses them, the $2l+1$ cancelling against the prefactor.  Certificate \vtag{11}
checks \eqref{eq:coefficients-coincide} in exact rational arithmetic for every
$l+m$ odd with $l \le 8$.  \emph{Exact}, range-limited --- a proof of that finite
statement, not evidence for it.

\subsection{The result}
\label{sec:odd-result}

One term survives,
\begin{equation}
  X_{lm}(q) = (-1)^m \frac{(l+m)!!}{(l-m-1)!!} \frac{j_l(2\pi q)}{2\pi q} ,
  \qquad l+m \text{ odd} ,
  \label{eq:odd-radial}
\end{equation}
and with \eqref{eq:mode-visibility} the mode visibility is
\begin{equation}
  V_{lm}(u,v) = i^m e^{im\psi} \, N_{lm} \,
  \frac{(l+m)!!}{(l-m-1)!!} \frac{j_l(2\pi q)}{q} ,
  \qquad l+m \text{ odd} ,
  \label{eq:odd-mode}
\end{equation}
the $(-i)^m$ and the $(-1)^m$ having combined into $i^m$.

At $q = 0$ the ratio in \eqref{eq:odd-radial} is indeterminate and takes its limit,
\begin{equation}
  \frac{j_l(a)}{a} \longrightarrow
  \begin{cases}
    \tfrac13 , & l = 1 , \\[2pt]
    0 , & l \ge 2 ,
  \end{cases}
  \label{eq:odd-limit}
\end{equation}
from $j_l(a) \approx a^l/(2l+1)!!$.  This is part of the statement rather than an
implementation detail, and certificate \vtag{13} covers it by putting $q = 0$ in
its grid.

As an independent check, \eqref{eq:odd-radial} agrees with equation (35) of
\citet{dholakia2025} once two conventions are accounted for: they use associated
Legendre functions without the Condon--Shortley phase, which contributes
$(-1)^m$, and they drop the harmonic normalisation, which \eqref{eq:odd-mode}
restores.

\subsection{What was checked}
\label{sec:odd-register}

\begin{center}
\begin{tabular}{lll}
\hline
tag & claim & rung \\
\hline
\vtag{11} & the collapse coefficients coincide, \eqref{eq:coefficients-coincide} & exact, $l \le 8$ \\
\vtag{13} & \eqref{eq:odd-radial} against the quadrature, $q = 0$ included & numerical, $l \le 5$ \\
\vtag{18} & the $K$ identity \eqref{eq:K-identity} & numerical, $l \le 6$ \\
\hline
\end{tabular}
\end{center}

The quadrature used by \vtag{13} is the one of
section~\ref{sec:visibility-register}, and it is never
dispatched to a closed form.  Were it otherwise, that certificate would compare the
formula with itself and could not fail.
